% ------------------------------------------------------------------
% CMPT 225 -- Assignment 1: Student Marking Rubric (one page, Letter)
% Compile with pdfLaTeX (Overleaf default). No extra files required.
% ------------------------------------------------------------------
\documentclass[10pt]{article}

% ---------- Page and fonts ----------
\usepackage[letterpaper,top=0.5in,bottom=0.5in,left=0.55in,right=0.55in]{geometry}
\usepackage[T1]{fontenc}
\usepackage[utf8]{inputenc}
\usepackage{lmodern}
\usepackage{microtype}
\usepackage{amsmath,amssymb}
\renewcommand{\familydefault}{\sfdefault}

% ---------- Tables, colour, lists ----------
\usepackage{array}
\usepackage{tabularx}
\usepackage[table]{xcolor}
\usepackage{enumitem}

\pagestyle{empty}
\setlength{\parindent}{0pt}
\setlength{\parskip}{0pt}
\setlength{\tabcolsep}{4pt}
\renewcommand{\arraystretch}{1.25}

\definecolor{hdr}{HTML}{2F3E4E}
\definecolor{tint}{HTML}{E8EDF2}
\definecolor{band}{HTML}{F5F7F9}
\definecolor{gridgray}{HTML}{9AA5B1}
\arrayrulecolor{gridgray}

% ---------- Helpers ----------
\newcolumntype{L}{>{\raggedright\arraybackslash}X}
\newcolumntype{C}{>{\centering\arraybackslash}X}
\newcommand{\hw}[1]{\textcolor{white}{\textbf{#1}}}
\newcommand{\hcell}[2]{\hw{#1}\newline{\footnotesize\textcolor{white}{#2}}}
\newcommand{\crit}[2]{\cellcolor{tint}\textbf{#1}\newline{\footnotesize\itshape #2}}
\newcommand{\sectionlabel}[1]{{\bfseries\color{hdr}#1}\par\vspace{2pt}}

\begin{document}

% ---------- Title ----------
\noindent{\LARGE\bfseries CMPT 225 Assignment 1 Rubric\par}
% ---------- Marking criteria ----------
\sectionlabel{}
{\small
\noindent\begin{tabularx}{\textwidth}{|>{\raggedright\arraybackslash}p{0.16\textwidth}|*{4}{L|}}
\hline
\rowcolor{hdr}
\hw{Criterion}
  & \hcell{Full credit}{complete and correct}
  & \hcell{Most credit}{minor flaws only}
  & \hcell{Partial credit}{significant gaps}
  & \hcell{Little or no credit}{missing or mostly wrong} \\
\hline

\crit{Correctness}{Applies to: all parts}
  & Every requested item is correct and clearly stated; nothing required is missing.
  & Correct in substance, with one small slip that does not change the main result.
  & Some items are correct, but a significant error or omission affects the result.
  & Mostly incorrect, off topic, or blank. \\
\hline

\crit{Reasoning and work shown}{Applies to: any justification or explanation}
  & A complete, logical argument; every step supports the answer; length and form match what is asked.
  & Sound reasoning with one small gap or skipped step that is easy to fill in.
  & Vague or partly incorrect, or the answer is restated without support.
  & No supporting work, or the work contradicts the answer. \\
\hline

\crit{Definitions, notation and conventions}{Applies to: A, C, D2--D3}
  & Definitions are precise and complete; symbols and stated conventions are used correctly.
  & Precise apart from one small omission or an inconsistent symbol.
  & The idea is recognizable but stated loosely, or conventions are applied inconsistently.
  & The definition is incorrect, or notation is misused so that the meaning is lost. \\
\hline

\crit{Complexity analysis}{Applies to: C1, D2--D3, E2, G1}
  & Uses the right variable, and is justified as required.
  & Correct bound, with a small flaw in the supporting argument.
  & The bound is valid but stated convention is ignored.
  & The bound is incorrect, unsupported, or missing. \\
\hline

\crit{Diagrams}{Applies to: D1(a), E1(a)--(b), E3(a)}
  & All required nodes and references drawn, correctly directed and labelled; null shown as $\varnothing$.
  & Structure is right apart from one missing or misdirected arrow, or a missing label.
  & Several references are wrong or missing, or the drawing is ambiguous.
  & The structure is fundamentally wrong, or not drawn as boxes and arrows. \\
\hline

\crit{Code}{Applies to: D3(a), E2, E3(c), G1(b)}
  & Correct in general and for each special case named; meets every stated constraint.
  & Correct overall, with one minor defect in a special case.
  & A reasonable idea, but it fails on some inputs or breaks a stated constraint.
  & The approach would not work, is unrelated to the question, or is missing. \\
\hline

\crit{Tracing and debugging}{Applies to: D1, D3(c), E3}
  & State is tracked accurately, step by step; a fault's cause and effect are explained.
  & One step is mis-tracked, but the rest of the trace is sound.
  & Tracing is partly correct, or a fault's symptoms are described without its cause.
  & State is not tracked, or the conclusions do not match the code given. \\
\hline

\crit{Clarity and presentation}{Applies to: all parts}
  & Answers are labelled by question number, legible and concise; code is indented and readable.
  & Occasional lapses that do not hinder reading.
  & Hard to follow in places; the marker must guess at your intent.
  & Illegible or disorganized, so credit cannot be assigned. \\
\hline
\end{tabularx}}
\vspace{9pt}

% ---------- Marking notes ----------
\sectionlabel{How marks are awarded}
{\small
\begin{itemize}[leftmargin=1.1em,itemsep=1pt,topsep=1pt,parsep=0pt,partopsep=0pt]
  \item Where a question splits its marks into labelled components (for example, a verdict and its justification, or running time and space), each component is assessed separately against the criteria above.
  \item A correct final answer without the requested justification or work earns only part of the marks.
  \item Code may be written in Java or in clear Java-like pseudocode. Use only the library classes a question explicitly allows.
\end{itemize}}
\vspace{8pt}

% ---------- Self-check ----------
\noindent\fcolorbox{gridgray}{band}{\parbox{\dimexpr\textwidth-2\fboxsep-2\fboxrule\relax}{\small\raggedright
\textbf{Before you submit:}\hspace{0.8em}%
\mbox{$\square$~every sub-part answered and labelled}\hspace{0.8em}%
\mbox{$\square$~work shown where requested}\hspace{0.8em}%
\mbox{$\square$~stated conventions followed}\hspace{0.8em}%
\mbox{$\square$~special cases checked in your code}\hspace{0.8em}%
\mbox{$\square$~name and student ID filled in}}}

\end{document}